\section{Installation, intégration et reproduction}

\subsection{Installer le prototype livré}

Chrome ne charge pas directement une archive ZIP. La procédure reproductible
est :

\begin{enumerate}
  \item décompresser \code{release/kora-extension.zip}, ou utiliser directement
        \code{release/kora-unpacked};
  \item ouvrir \code{chrome://extensions};
  \item activer \textbf{Mode développeur};
  \item cliquer \textbf{Charger l’extension non empaquetée};
  \item sélectionner le dossier qui contient directement
        \code{manifest.json};
  \item ouvrir ou actualiser la page financière;
  \item épingler KORA, ouvrir sa popup et choisir \textbf{Fulfulde}.
\end{enumerate}

\begin{keybox}[Démonstration minimale en moins d’une minute]
Vérifier d’abord un montant et un compte sur la page source. Activer Fulfulde,
constater que les libellés changent mais que les valeurs restent identiques.
Changer d’étape, provoquer une erreur simulée, puis choisir Français ou
English pour restaurer la page.
\end{keybox}

\subsection{Construire depuis les sources}

\textbf{Prérequis :} Node.js 20 ou supérieur, npm, Python 3 et Google Chrome.
PyTorch et NumPy sont nécessaires uniquement pour réentraîner le réseau; les
poids versionnés suffisent pour construire et tester l’extension.

\begin{lstlisting}[language=bash]
git clone https://github.com/kaezerwatto/KORA.git
cd KORA
npm install
npm run build
npm test
npm run package
\end{lstlisting}

\code{npm run package} reconstruit la terminologie, regroupe le bundle,
exporte le modèle JSON, produit le dossier unpacked et crée le ZIP. Le script
interrompt la livraison au-delà de 10 Mio.

\subsection{Régénérer les ressources linguistiques}

\begin{lstlisting}[language=bash]
npm run build:corpora       # inventaire et empreintes des deux volumes
npm run build:terminology   # index FR/EN -> fub et codes proteges
npm run build:model         # modele n-gramme
npm run model:train         # reseau, early stopping, int8
npm run model:export        # JSON portable + empreinte SHA-256
npm run model:demo          # predictions et cas d'abstention
\end{lstlisting}

L’entraînement écrit \code{src/fub-neural-model.js}; le build l’embarque dans
\code{content.js}. L’export indépendant est
\code{release/kora-fub-neural-model.json}. Il contient architecture,
vocabulaire, tenseurs int8 encodés, échelles et empreinte de contrôle, mais
aucun contenu utilisateur.

\subsection{Utiliser le moteur complet séparément}

Dans un projet JavaScript ESM, importer \code{translateText}. Cette API porte
la protection des valeurs et convient à une traduction contrôlée :

\begin{lstlisting}[language={}]
import { translateText } from "./src/engine.js";

const result = translateText("sodle insufisant", {
  sourceHint: "fr",
  domainHint: "financial",
  protectedValues: []
});

console.log(result.text);    // Kaalisi yonaa
console.log(result.method);  // spelling-memory
\end{lstlisting}

\begin{tabularx}{\linewidth}{@{}L{38mm}Y@{}}
\toprule
\textbf{Option/résultat} & \textbf{Signification}\\
\midrule
\code{sourceHint} & \code{auto}, \code{fr} ou \code{en}.\\
\code{domainHint} & \code{general} ou \code{financial}.\\
\code{protectedValues} & Chaînes à conserver strictement.\\
\code{text} & Traduction acceptée ou texte original.\\
\code{changed} & Indique si le rendu a changé.\\
\code{method} &
\code{memory}, \code{spelling-memory}, \code{template},
\code{glossary}, \code{protected}, \code{integrity-fallback} ou
\code{original}.\\
\code{fluency} & Score du modèle lorsque pertinent.\\
\bottomrule
\end{tabularx}

\subsection{Utiliser uniquement les modèles de langue}

\begin{lstlisting}[language={}]
import {
  predictNextFub,
  rankFubCandidates,
  scoreFubFluency
} from "./src/language-model.js";

predictNextFub("transfeer maa", 5);
scoreFubFluency("Kaalisi yonaa");
rankFubCandidates([
  "Transfeer maa neldaama.",
  "Transfeer maa yettaama."
]);
\end{lstlisting}

L’API expose trois fonctions : \code{predictNextFub} propose des mots
suivants, \code{scoreFubFluency} évalue une séquence \code{fub} et
\code{rankFubCandidates} classe des formulations. Elles ne connaissent ni
la source française ou anglaise, ni les règles financières : elles ne doivent
pas être utilisées comme traducteur autonome.

Pour un runtime non JavaScript, le JSON exporté permet de reproduire
l’inférence : décoder les tenseurs int8, récupérer les embeddings des trois
mots, calculer la couche tanh, projeter sur le vocabulaire et appliquer le
softmax. \code{src/neural-language-model.js} est l’implémentation de référence.

\subsection{Ajouter une langue ou un domaine}

Pour une autre langue, l’architecture reste réutilisable. Il faut cependant
plus qu’un simple remplacement de chaînes :

\begin{enumerate}
  \item définir le code, l’écriture et les variantes cibles;
  \item constituer une mémoire source$\rightarrow$cible avec provenance;
  \item adapter les patrons morphologiques et la détection;
  \item entraîner ou remplacer le modèle de fluidité;
  \item créer des tests d’intégrité et de parité FR/EN;
  \item faire relire les actions critiques et alertes;
  \item versionner les décisions et publier les limites.
\end{enumerate}

\subsection{Dépendances et services déclarés}

\begin{longtable}{@{}L{43mm}L{35mm}L{50mm}L{28mm}@{}}
\toprule
\textbf{Composant} & \textbf{Version/forme} & \textbf{Rôle} &
\textbf{Dans le ZIP}\\
\midrule
\endfirsthead
\toprule
\textbf{Composant} & \textbf{Version/forme} & \textbf{Rôle} &
\textbf{Dans le ZIP}\\
\midrule
\endhead
esbuild & 0.28.2 & Bundle local du content script & Non\\
puppeteer-core & 25.17.0 & Tests navigateur et capture & Non\\
Node test runner & Node.js & Tests unitaires & Non\\
PyTorch + NumPy & Environnement Python optionnel & Réentraînement & Non\\
Chrome MV3 & API native & Exécution de l’extension & Oui, contrat\\
ElevenLabs & API de production vidéo & Narration du pitch public & Non\\
Web Audio API & API navigateur & Musique originale du montage & Non\\
Modèles KORA & JavaScript int8 + n-gramme & Fluidité locale & Oui\\
\bottomrule
\end{longtable}

ElevenLabs n’appartient pas au runtime KORA. Seul le script public de la vidéo
est transmis lors de la production; aucune page, saisie ou donnée financière
n’est envoyée. Les clés restent dans un fichier local non distribué.

\subsection{Limites connues et évolution}

\begin{tabularx}{\linewidth}{@{}L{48mm}Y@{}}
\toprule
\textbf{Limite} & \textbf{Traitement / prochaine étape}\\
\midrule
Phrase libre totalement inconnue &
Conserver la source; enrichir ensuite par une entrée sourcée et revue.\\
Terminologie candidate &
Double revue à l’aveugle par locuteurs \code{fub} et test de compréhension.\\
Variations régionales &
Échantillonner Ngaoundéré, Garoua et Maroua; versionner les variantes.\\
Image, canvas, Shadow DOM fermé &
Hors accès standard; envisager un module OCR séparé avec consentement.\\
Certification bancaire &
Audit externe, gouvernance des versions, essais sécurité et conformité.\\
Déploiement navigateur &
Signature, politique de mises à jour et publication sur un canal contrôlé.\\
\bottomrule
\end{tabularx}

\begin{warningbox}[Priorité de feuille de route]
La prochaine amélioration n’est pas d’augmenter la génération libre, mais de
faire valider davantage de concepts et de mesurer le taux d’abstention sur un
jeu inédit. Cette approche augmente la surface utile sans dégrader la sûreté
financière.
\end{warningbox}
